\section{Conclusion}
\label{sec:conclusion}

A property survives representation change when the inference that defines it
survives.  That requires more than the persistence of a symbol, numerical
agreement on selected examples, or an assertion that two descriptions encode
the same physics.  It requires a typed bearer, a physical probe class,
property-specific fidelity, explicit loss, theorem-backed recovery, and the
uniformity needed by the limiting claim.

The uniform recovery-to-spectrum theorem makes this criterion concrete.
Positive Euclidean semigroup bounds with a common physical rate imply a target
spectral exclusion only after their quadratic forms are recovered on a dense
centered core.  The hidden-sector, collapsing-gap, and noncommuting-limit
examples show that none of those qualifiers is cosmetic.  The Yang--Mills
case then demonstrates the broader lesson: locality, gauge invariance,
positivity, scale flow, normalization, completeness, and spectral meaning are
distributed across descriptions, so infrared rigidity must be generated and
transported through a chain of distinct theorems.

This supports a restrained big picture.  Physical objectivity may be
understood as certified coherence across partial descriptions, while the
choice between structural-realist, entity-realist, empiricist, and
perspectival ontologies remains open.  The immediate mathematical task is
sharper: construct a noncircular, full-measure rigidity estimate uniform in
volume and cutoff; build the continuum OS-positive Yang--Mills theory; and
prove that the controlled observable core is physically complete.  Until
those certificates close, a regulated gap is a property of a representation,
not yet the mass gap of the continuum theory.
